\documentclass[10pt, letterpaper]{article}

% Packages:
\usepackage[
    ignoreheadfoot,
    top=2 cm,
    bottom=2 cm,
    left=2 cm,
    right=2 cm,
    footskip=1.0 cm,
]{geometry}
\usepackage[explicit]{titlesec}
\usepackage{tabularx}
\usepackage{array}
\usepackage[dvipsnames]{xcolor}
\definecolor{primaryColor}{RGB}{0, 79, 144} % blue
\usepackage{enumitem}
\usepackage{fontawesome5}
\usepackage{amsmath}
\usepackage[
    pdftitle={Chunxu Han's CV},
    pdfauthor={Chunxu Han},
    pdfcreator={LaTeX with RenderCV},
    colorlinks=true,
    urlcolor=primaryColor
]{hyperref}
\usepackage[pscoord]{eso-pic}
\usepackage{calc}
\usepackage{bookmark}
\usepackage{lastpage}
\usepackage{changepage}
\usepackage{paracol}
\usepackage{ifthen}
\usepackage{needspace}
\usepackage{iftex}
\usepackage{multicol}
\usepackage{datetime2}
\usepackage{graphicx}

\ifPDFTeX
    \input{glyphtounicode}
    \pdfgentounicode=1
    \usepackage[T1]{fontenc}
    \usepackage[utf8]{inputenc}
    \usepackage{lmodern}
\fi

\usepackage[default, type1]{sourcesanspro}

\AtBeginEnvironment{adjustwidth}{\partopsep0pt}
\pagestyle{empty}
\setcounter{secnumdepth}{0}
\setlength{\parindent}{0pt}
\setlength{\topskip}{0pt}
\setlength{\columnsep}{0.15cm}
\makeatletter
\let\ps@customFooterStyle\ps@plain
\patchcmd{\ps@customFooterStyle}{\thepage}{
    \color{gray}\textit{\small Chunxu Han - Page \thepage{} of \pageref*{LastPage}}
}{}{}
\makeatother
\pagestyle{customFooterStyle}

\titleformat{\section}{
    \needspace{4\baselineskip}
    \Large\color{primaryColor}
}{}{}{
    \textbf{#1}\hspace{0.15cm}\titlerule[0.8pt]\hspace{-0.1cm}
}[]

\titlespacing{\section}{-1pt}{0.3 cm}{0.2 cm}

\newenvironment{highlights}{
    \begin{itemize}[
        topsep=0.10 cm,
        parsep=0.10 cm,
        partopsep=0pt,
        itemsep=0pt,
        leftmargin=0.4 cm + 10pt
    ]
}{
    \end{itemize}
}

\newenvironment{highlightsforbulletentries}{
    \begin{itemize}[
        topsep=0.10 cm,
        parsep=0.10 cm,
        partopsep=0pt,
        itemsep=0pt,
        leftmargin=10pt
    ]
}{
    \end{itemize}
}

\newenvironment{onecolentry}{
    \begin{adjustwidth}{0.2 cm + 0.00001 cm}{0.2 cm + 0.00001 cm}
}{
    \end{adjustwidth}
}

\newenvironment{twocolentry}[2][]{
    \onecolentry
    \def\secondColumn{#2}
    \setcolumnwidth{\fill, 4.5 cm}
    \begin{paracol}{2}
}{
    \switchcolumn \raggedleft \secondColumn
    \end{paracol}
    \endonecolentry
}

\newenvironment{threecolentry}[3][]{
    \onecolentry
    \def\thirdColumn{#3}
    \setcolumnwidth{1 cm, \fill, 4.5 cm}
    \begin{paracol}{3}
    {\raggedright #2} \switchcolumn
}{
    \switchcolumn \raggedleft \thirdColumn
    \end{paracol}
    \endonecolentry
}

\newenvironment{header}{
    \setlength{\topsep}{0pt}\par\kern\topsep\centering\color{primaryColor}\linespread{1.5}
}{
    \par\kern\topsep
}

\let\hrefWithoutArrow\href
\renewcommand{\href}[2]{\hrefWithoutArrow{#1}{\ifthenelse{\equal{#2}{}}{ }{#2 }\raisebox{.15ex}{\footnotesize \faExternalLink*}}}


\begin{document}
    \newcommand{\AND}{\unskip
        \cleaders\copy\ANDbox\hskip\wd\ANDbox
        \ignorespaces
    }
    \newsavebox\ANDbox
    \sbox\ANDbox{}

    \begin{header}
        \begin{minipage}[t]{0.78\linewidth}
        \raggedright
        \fontsize{28 pt}{28 pt}
        \textbf{Chunxu Han}

        \vspace{0.3 cm}

        \normalsize
        \mbox{{\footnotesize\faMapMarker*}\hspace*{0.13cm}Peder Skrams Gade 16A, København K, Denmark}%
        \kern 0.25 cm%
        \AND%
        \kern 0.25 cm%
        \mbox{\hrefWithoutArrow{mailto:s220311@dtu.dk}{{\footnotesize\faEnvelope[regular]}\hspace*{0.13cm}s220311@dtu.dk}}%
        \kern 0.25 cm%
        \AND%
        \kern 0.25 cm%
        \mbox{\hrefWithoutArrow{tel:+45 52 73 06 85}{{\footnotesize\faPhone*}\hspace*{0.13cm}+45 52 73 06 85}}%
        \kern 0.25 cm%
        \AND%
        \kern 0.25 cm%
        \mbox{\hrefWithoutArrow{https://www.linkedin.com/in/chunxu-han/}{{\footnotesize\faLinkedinIn}\hspace*{0.13cm}chunxu-han}}%
        \kern 0.25 cm%
        \AND%
        \kern 0.25 cm%
        \mbox{\hrefWithoutArrow{https://github.com/chunxubioinfor}{{\footnotesize\faGithub}\hspace*{0.13cm}chunxubioinfor}}%
        \end{minipage}%
        \hfill%
        \begin{minipage}[t]{0.18\linewidth}
        \raggedleft
        \vspace{-1.2cm}
        \includegraphics[width=2.4cm, height=2.8cm, keepaspectratio=false]{portrait.jpg}
        \end{minipage}
    \end{header}

    \vspace{0.6 cm}

    \begin{onecolentry}
        \textit{While my background is rooted in bioinformatics, I have made a deliberate effort to build familiarity with clinical research and its data environment. Through the Synapse Clinical Winter School, I gained exposure to the fundamentals of clinical trials and developed an understanding of central domain concepts and terminology used in clinical development. That experience made me even more interested in applying my quantitative training and strong R programming skills in a setting where data quality and analysis contribute to decisions around safety, efficacy, and treatment development. I would be excited to continue growing in this area and contribute to rigorous, clinically relevant data workflows.}
    \end{onecolentry}

    \vspace{0.3 cm}

    \section{Relevant Experience}

    \begin{onecolentry}

        \textbf{R Programming and Software Development} \\
        R is my primary language and I have used it in production settings for two years. At DTU HealthTech, I contributed to \href{https://github.com/kvittingseerup/IsoformSwitchAnalyzeR}{IsoformSwitchAnalyzeR} v2, an R package used by researchers worldwide for isoform-level transcriptomic analysis. My work covered feature development, bug fixing, triaging GitHub issues from external users, and writing reproducible analysis scripts. I understand what it means to write R code that others will run and rely on, where correctness and documentation matter as much as functionality.

        \vspace{0.3 cm}

        \textbf{Quantitative Analysis and Data Quality} \\
        At Novo Nordisk, my thesis centred on developing a novel metric, the Batch Mixing Index, to detect and quantify systematic variation in large single-cell manufacturing datasets. I designed the statistical formulation, implemented it in Python, and validated it using permutation-based significance testing. Across both roles, I have worked extensively with complex, heterogeneous biological datasets where data quality checks, consistency validation, and careful documentation were essential parts of the workflow. I take accuracy seriously and have developed habits around QC that I carry into every analysis.

        \vspace{0.3 cm}

        \textbf{Clinical and Pharma Exposure} \\
        Through the Synapse Clinical Winter School, I gained foundational knowledge of clinical trial design, phases of drug development, and the data environment of clinical research. My consulting internship at Illuminera, IQVIA gave me direct exposure to the pharmaceutical industry, working on client-facing projects and collaborating with cross-functional teams across data management and strategy. I also have a GMP compliance background from a QC internship at Sinovac Biotech and relevant coursework, which has given me an appreciation for the regulatory standards and documentation rigour that underpin clinical data work.

    \end{onecolentry}

    %% ===== FIXED: Do not edit below this line =====

    \section{Education}

        \begin{twocolentry}{
    Lyngby, Denmark

    Sept 2023 -- March 2026
}
    \textbf{MSc in Bioinformatics and Systems Biology}, Technical University of Denmark
    \begin{highlights}
        \item \textbf{Specialisation}: Biomedical Bioinformatics
        \item Enrolled in both \textbf{Machine Learning} and \textbf{Deep Learning} courses, enabling me to develop an \textbf{AI-driven mindset} when formulating projects and developing new bioinformatics pipelines.
    \end{highlights}
\end{twocolentry}

        \vspace{0.3 cm}

\begin{twocolentry}{
    Beijing, China

    Sept 2017 -- June 2021
}
    \textbf{BSc in Biotechnology}, Beijing Normal University
    \begin{highlights}
        \item \textbf{Coursework:} Built a solid foundation in \textbf{Biotech} with courses covering ecology, molecular biology, cell biology, and bioinformatics, alongside \textbf{lab training}.
    \end{highlights}
\end{twocolentry}


    \section{Work Experience}

        \begin{twocolentry}{
    Målev, Denmark

    Sep 2025 -- March 2026
}
    \textbf{Master Student}, Novo Nordisk

    \vspace{0.1cm}

    My master's thesis focuses on developing computational frameworks to evaluate batch-to-batch variability in single-cell RNA-seq data:

    \begin{highlights}
        \item \textbf{Methodology Development:} Identified limitations in existing batch effect methods and developed the entropy-based Batch Mixing Index (BMI) to quantify per-cell neighbourhood mixing with statistical testing. This work strengthened my ability to translate methodological gaps into implementable solutions.

        \vspace{0.1cm}

        \item \textbf{Meetings and Collaboration:} Participated in cross-functional meetings and collaborated with R\&D and CMC teams to ensure biological interpretability and industrial relevance of project.
    \end{highlights}
\end{twocolentry}

        \vspace{0.3 cm}

        \begin{twocolentry}{
    Lyngby, Denmark

    Oct 2023 -- July 2025
}
    \textbf{Bioinformatics Research Assistant}, DTU HealthTech

    \vspace{0.1cm}

    I worked in Isoform Analysis Group for two years, which has been a transformative experience. I have grown from a student who only focuses on theoretical studies into a bioinformatician to deal with real-world problems:

    \begin{highlights}
        \item \textbf{Development and Maintenance of Tools:} Contributed to developing and updating the \href{https://github.com/kvittingseerup/IsoformSwitchAnalyzeR}{IsoformSwitchAnalyzeR} R package. This experience sharpened my coding skills and can-do mindset, especially when handling issues with tools I'd never used before. I pushed myself to attempt and now I'm comfortable with the steep learning curve.

        \vspace{0.1cm}

        \item \textbf{Bioinformatics Pipelines:} Developed and optimized bioinformatics pipelines, supporting others in running analyses. I mastered multi-tasking, learned to leverage \textbf{HPC} for parallel processing, and \textbf{Git} to achieve version control.
    \end{highlights}
\end{twocolentry}

        \vspace{0.3 cm}

        \begin{twocolentry}{
    Shanghai, China

    Feb 2023 -- May 2023
}
    \textbf{Consulting Intern}, Illuminera, IQVIA

    \vspace{0.1cm}

    My first experience in a business setting, where I gained insights into the pharmaceutical market and developed professional skills:

    \begin{highlights}
        \item \textbf{Qualitative and Quantitative Research:} Conducted market trend and competitor analysis to formulate actionable strategies to clients. I also gained a deeper understanding of the pharmaceutical market in China and globally.

        \vspace{0.1cm}

        \item \textbf{Client and KOLs Engagement:} I learnt how to build professional relationships with clients and KOLs and coordinate with stakeholders to align project progress.

        \vspace{0.1cm}

        \item \textbf{Data-Driven Deliverables:} Utilized tools like PowerPoint and Power BI to deliver strategic reports, transforming complex data into actionable business insights.
    \end{highlights}
\end{twocolentry}

        \vspace{0.3 cm}



    \section{Skills and Hobbies}

        \begin{onecolentry}
            \textbf{Programming:} Python, R, SAS, Bash, SQL, Snakemake
        \end{onecolentry}

        \vspace{0.2 cm}

        \begin{onecolentry}
            \textbf{Languages:} English (proficient), Chinese (native), Danish (Basic)
        \end{onecolentry}

        \vspace{0.2 cm}

        \begin{onecolentry}
            \textbf{Soft Skills:} Communication, Teamwork, Client Engagement, Multitasking, Critical Thinking, Strong Work Ethic
        \end{onecolentry}

        \vspace{0.2 cm}

        \begin{onecolentry}
            \textbf{Hobbies:}
            I've been playing \textbf{badminton} since an early age. It's a hobby that keeps me focused and energized. During holidays, I enjoy \textbf{traveling} around Europe to experience new food and cultures. I also love \textbf{hiking}, both for the chance to connect with nature and for the challenge of seeing hidden scenery that isn't easily accessible. After moving to Denmark, I discovered a new passion: \textbf{sailing}. Last summer, we sailed to Ven and Hallands Väderö, and it was a completely new experience and I feel free to ride the infinite waves.
        \end{onecolentry}


\end{document}
